\chapter{Resultados}
\label{cap:resultados}

\section{Equivalencia entre hardware y simulación}

El resultado más fuerte del proyecto es también el más simple de enunciar:
\textbf{el port reproduce el diseño original bit a bit}. Con el flujo
\cmd{tcela-17} decodificado completo en la placa y en la simulación del
núcleo, las dos memorias de cuadros resultaron idénticas byte a byte, en los
tres planos y en los cuatro buffers. La misma igualdad se observó en los
valores de error contra la referencia de todos los flujos comparados por
ambos caminos, idénticos hasta el tercer decimal.

Esto delimita responsabilidades (sección~\ref{sec:limites-simulacion}):
cualquier diferencia que quede contra el decodificador de referencia ya
estaba en el diseño original, y no fue introducida por el port.

\section{Precisión de la IDCT}

La IDCT del núcleo cumple las seis condiciones del estándar IEEE~1180-1990
(tabla~\ref{tab:ieee1180}), con el error pico en el máximo permitido de~1 y
los errores cuadráticos y medios uno o dos órdenes de magnitud por debajo de
sus límites.

\begin{table}[ht]
  \centering\small
  \begin{tabular}{@{}l r r r@{}}
    \toprule
    \textbf{Métrica} & \textbf{Peor caso medido} & \textbf{Límite} & \\
    \midrule
    Error pico                       & 1        & 1      & cumple \\
    Error cuadrático por posición    & 0,0052   & 0,06   & cumple \\
    Error cuadrático global          & 0,0036   & 0,02   & cumple \\
    Error medio por posición         & 0,0016   & 0,015  & cumple \\
    Error medio global               & 0,000056 & 0,0015 & cumple \\
    \bottomrule
  \end{tabular}
  \caption{IDCT contra IEEE 1180-1990: peor valor entre las seis
    condiciones.}
  \label{tab:ieee1180}
\end{table}

\section{Conformidad}
\label{sec:res-conformidad}

\paragraph{Imágenes intra}
En las imágenes codificadas sin predicción, la salida del hardware difiere
de la referencia en un MAD de 0,004 a 0,005 con error pico~1, en luma y en
crominancia, en cuatro flujos distintos. Es la única diferencia que el
estándar admite entre dos implementaciones conformes de la IDCT. Un flujo
completo, \cmd{mcp10ccett} ($720 \times 576$, con imágenes I, P y B), es
bit-exacto contra la referencia en todos sus buffers y planos.

\paragraph{Batería de conformidad}
La tabla~\ref{tab:conformidad} resume los veredictos sobre los 39 flujos;
el detalle está en el anexo~\ref{anx:regresion}.

\begin{table}[ht]
  \centering\small
  \begin{tabularx}{\linewidth}{@{}l r X@{}}
    \toprule
    \textbf{Veredicto} & \textbf{Flujos} & \textbf{Observación} \\
    \midrule
    Bit-exacto      & 3  & \cmd{mcp10ccett}, \cmd{tcela-15}, \cmd{tcela-17} \\
    Cercano         & 34 & MAD $\leq 0{,}04$ y pico $\leq 4$ en el peor buffer \\
    Deriva          & 0  & \\
    Sin escritura   & 2  & MPEG-1, que el núcleo no decodifica \\
    \bottomrule
  \end{tabularx}
  \caption{Veredictos de conformidad (39 flujos, hasta 8 imágenes cada uno).}
  \label{tab:conformidad}
\end{table}

Los 37 flujos MPEG-2 decodifican bit-exactos o cercanos, y en los cercanos
la diferencia es la que admite el estándar entre dos IDCT conformes: el
peor buffer de toda la batería tiene MAD 0,038 (\cmd{att-mismatch}, un
flujo de $32 \times 32$ hecho justamente para estresar la IDCT), y en 30 de
los 34 el peor buffer está por debajo de 0,02. Ningún flujo deriva. Para
llegar ahí hubo que corregir la referencia
(sección~\ref{sec:conformidad-referencia}), la forma de cortar los flujos y
dos errores del diseño original (sección~\ref{sec:conformidad-upstream}).

Los dos flujos MPEG-1 (\cmd{mei-16-long}, \cmd{mei-16v2}) no tienen
\emph{sequence extension}: el núcleo es solo MPEG-2, y el propio diseño
original los marca como no soportados.

\cmd{tcela-15} figuraba también como \emph{sin escritura}. Es un flujo de
prueba de relleno: cada imagen trae unos pocos cientos de bytes de datos y
después unos 60~KB de ceros. Su contenido es gris medio uniforme, que la
memoria de cuadros guarda como ceros, y la comparación consideraba
\enquote{nunca escrito} a todo buffer uniforme. Ahora un buffer se considera
vacío solo si tiene el valor con que el núcleo limpia la memoria al
arrancar o el que la captura escribe antes de cada prueba. Con ese
criterio, \cmd{tcela-15} es bit-exacto. Por el servicio de red también
entrega sus 8 imágenes completas.

\subsection{Un error en la referencia}
\label{sec:conformidad-referencia}

La primera versión de esta batería daba un resultado muy distinto: un solo
flujo bit-exacto, 12 cercanos y 23 con deriva, de los cuales once fallaban
desde la primera imagen con MAD de 12 a 18. Todos ellos cargan una matriz de
cuantización propia en la cabecera de secuencia, y casi todos cargan
literalmente la matriz por defecto, lo que debería ser equivalente a no
cargar ninguna. La conclusión natural fue un error en el camino de carga de
matrices del núcleo, y así figuró como limitación conocida.

Dos observaciones la desmintieron. La documentación del diseño original
declara aprobados todos esos flujos. Y una simulación del núcleo con un
monitor sobre la memoria de la matriz mostró que las 64 escrituras llegaban
a la dirección correcta con el valor correcto. Al mirar las imágenes en
lugar de los números, la del hardware era nítida y la de la referencia una
imagen en bloques. El error estaba en la copia de \cmd{mpeg2decode} que
acompaña al núcleo: en la cabecera de secuencia, la lectura de las matrices
había quedado dentro de un \cmd{if} que solo se ejecuta con las trazas
activadas. Sin trazas, que es como la usa la comparación, la referencia no
cargaba la matriz ni consumía sus 512~bits.

Corregido ese \cmd{if}, las capturas del hardware, idénticas byte a byte a
las anteriores, pasaron a 2 bit-exactos, 29 cercanos y 5 con deriva:
\cmd{gi4} bajó de MAD 17,9 a 0,005 en su imagen intra. La corrección también resolvió la mayor parte de la deriva en
imágenes predichas, porque varios de esos flujos cargan matrices para
bloques no intra.

\subsection{Dos errores en el diseño original}
\label{sec:conformidad-upstream}

Corregida la referencia quedaban cinco flujos con deriva. En dos de ellos
(\cmd{sony-ct1}, \cmd{hhi-burst-long}) el corte de 8 imágenes terminaba en
el primer campo de un cuadro; ahora el corte nunca separa los dos campos.
Los otros tres (\cmd{nokia6-60}, \cmd{tcela-10}, \cmd{tcela-14}) eran
errores del núcleo: la simulación del RTL original da el mismo volcado, bit
a bit, que la placa, así que no los introdujo el port, aunque la
documentación original da esos flujos por aprobados.

Simular hasta las imágenes afectadas lleva horas por volcado. El método que
funcionó fue otro: modificar una copia de la referencia para que se
comporte como se sospecha que se comporta el hardware, y ver si los
volcados de la placa coinciden con ella al nivel de ruido de la IDCT. Para
desempatar entre el hardware y la referencia se usó además un tercer
decodificador independiente (ffmpeg).

\begin{itemize}
  \item \textbf{Redondeo del promedio de predicciones.} Al promediar la
    predicción hacia adelante y hacia atrás, el núcleo redondeaba según el
    signo del píxel (que guarda con un desplazamiento de $-128$), en lugar
    de $(a+b+1)/2$ como pide la norma. Los píxeles oscuros perdían un nivel.
    En imágenes B el efecto no se acumula (MAD $\approx 0{,}26$), pero la
    predicción \emph{dual prime} usa ese promedio en imágenes P, que son
    referencia de las siguientes: en \cmd{nokia6-60} y \cmd{tcela-14} el
    sesgo crecía imagen a imagen hasta MAD~2. Con la referencia emulando
    ese redondeo, el hardware coincidía en MAD 0,02.
  \item \textbf{Orden de las matrices recargadas.} Las matrices de
    cuantización se transmiten siempre en orden zigzag, pero el núcleo las
    ubicaba con el orden de recorrido de la imagen anterior. Una segunda
    cabecera de secuencia que llega después de una imagen con recorrido
    alternativo, como en \cmd{tcela-10}, dejaba los pesos permutados. La
    firma estaba en el dominio de la DCT: la ganancia entre hardware y
    referencia era exactamente 1 en muchos coeficientes y, en los demás, un
    cociente de dos valores de la matriz (22/19, 16/19, 19/16). El
    comentario del código original ya decía que debía usarse el orden zigzag.
\end{itemize}

Los dos arreglos suman tres líneas de RTL. Con ellos, \cmd{nokia6-60} bajó
de MAD 1,79 a 0,026, \cmd{tcela-14} de 1,62 a 0,025 y \cmd{tcela-10} de
2,40 a 0,009, y las imágenes B del resto de la batería, de 0,26 a menos de
0,02.

\section{Determinismo y robustez}

Las diez decodificaciones de la prueba de determinismo, con y sin
\emph{reset}, dieron imágenes idénticas. En la prueba de robustez, ninguna de
las tres entradas dañadas dejó al decodificador trabado: el flujo correcto
siguiente se decodificó bit-exacto en los tres casos, solo con el cambio de
flujo y sin \emph{reset}. El indicador de error del núcleo se activó con el
flujo truncado y con los bits invertidos, y no con los bytes aleatorios, que
no contienen ningún código de inicio que el analizador pueda rechazar. El
\emph{watchdog} no se activó en ningún caso.

\section{Rendimiento}
\label{sec:res-rendimiento}

\begin{table}[ht]
  \centering\small
  \begin{tabular}{@{}l c r r r@{}}
    \toprule
    \textbf{Flujo} & \textbf{Resolución} & \textbf{img/s} & \textbf{desvío} & \textbf{$\times$ tiempo real} \\
    \midrule
    tek60          & $704 \times 480$ & 23,6 & 0,5 & 0,79 \\
    tek-5-long     & $704 \times 480$ & 22,7 & 0,7 & 0,76 \\
    mei-2-60f      & $704 \times 480$ & 21,2 & 1,0 & 0,71 \\
    tcela-7        & $720 \times 480$ & 19,9 & 0,9 & 0,66 \\
    nokia6-60      & $720 \times 576$ & 11,9 & 0,2 & 0,48 \\
    hhi-burst-long & $720 \times 576$ & 13,1 & 0,5 & 0,52 \\
    \midrule
    \emph{Tears of Steel} & $352 \times 288$ & 53,7 & 0,2 & 2,24 \\
    \bottomrule
  \end{tabular}
  \caption{Tasa de decodificación en la placa con el bitstream
    \cmd{0.1.0+918dde0} (5 corridas por flujo; 3 para el clip).}
  \label{tab:rendimiento}
\end{table}

La tabla~\ref{tab:rendimiento} muestra la tasa de decodificación. En
definición estándar el decodificador alcanza 20 a 23,5 imágenes por segundo en
$704/720 \times 480$ y 12 a 13 en $720 \times 576$, es decir, entre la
mitad y cuatro quintos del tiempo real. El límite está en la etapa de
reconstrucción del núcleo original (sección~\ref{sec:optimizacion}). En el
formato CIF del clip de demostración, $352 \times 288$, decodifica a 53,7
imágenes por segundo, 2,2 veces el tiempo real.

De punta a punta, con el cliente en otra computadora de la red local que
recibe cada imagen a través del servicio y la biblioteca, el clip llega a
unas 50--51 imágenes por segundo, idéntico byte a byte entre corridas.
Contra la referencia, el clip completo da un MAD medio de 0,010 (máximo
0,018 en un cuadro) con error pico de 3 sobre 255; antes de corregir los
errores de la sección~\ref{sec:conformidad-upstream} daba 0,073.

\section{Consumo de potencia}

La estimación de SmartPower para el diseño completo es de \textbf{1,66~W}
en condiciones típicas (1,96~W en el peor caso), de los cuales el MSS, con
sus cinco procesadores y la interfaz de memoria DDR4, explica el 90~\%. El
\textbf{periférico decodificador completo consume unos 68~mW}, el 4,3~\% del
total: 63,6~mW el núcleo y apenas 4~mW los tres bloques de integración. Dentro
del núcleo, la IDCT y el PLL propio son los mayores consumidores, con un
cuarto cada uno. El desglose está en el anexo~\ref{anx:potencia}.

Dos salvedades. Es una estimación sin vectores, no una medición. Y se hizo
sobre la netlist previa a la última síntesis, que agregó el DMA por bloques,
el identificador de versión y la interrupción de imagen: cambios menores,
que no deberían alterar el orden de magnitud.

\section{Demostración pública}

La figura~\ref{fig:demo-comparacion} muestra una imagen del clip decodificada
por la FPGA junto a la del decodificador de referencia y su diferencia
amplificada: la diferencia es invisible sin amplificar.

\begin{figure}[ht]
  \centering
  \includegraphics[width=\linewidth]{demo_comparacion_referencia}
  \caption{Salida de la FPGA (izquierda), decodificador de referencia
    \cmd{mpeg2decode} (centro) y diferencia absoluta amplificada 32 veces
    (derecha). Contenido: \emph{Tears of Steel} \citep{Blender_2012_tears}.}
  \label{fig:demo-comparacion}
\end{figure}

\begin{table}[ht]
  \centering\small
  \begin{tabular}{@{}l r r r@{}}
    \toprule
    \textbf{Condición} & \textbf{Placa} & \textbf{Llegada} & \textbf{Buffer inicial} \\
    \midrule
    JPEG, red a 2,4~Mbit/s (primera versión) & 2,11$\times$ & 0,60$\times$ & 15~s \\
    H.264, red a 2,4~Mbit/s                  & 2,09$\times$ & 2,00$\times$ & 1,5~s \\
    H.264, red a 1~Mbit/s                    & 2,06$\times$ & 1,12$\times$ & 1,5~s \\
    H.264, por el túnel público              & 2,02--2,05$\times$ & 1,98--2,01$\times$ & 2,3~s \\
    \bottomrule
  \end{tabular}
  \caption{Demostración en vivo: tasa de decodificación en la placa y de
    llegada al navegador, en múltiplos del tiempo real. Sin cortes durante la
    reproducción en ningún caso.}
  \label{tab:demo}
\end{table}

En la demostración en vivo (tabla~\ref{tab:demo}), la placa decodifica a
unas dos veces el tiempo real en todas las condiciones, porque ya no la
frena la red del visitante. La tasa de llegada depende de la red, y la
página calcula cuánto buffer necesita para reproducir sin cortes. Desde una
red externa real, el clip llegó a 2,06 veces el tiempo real.

\section{Síntesis}

\begin{itemize}
  \item El port es fiel: hardware y simulación coinciden bit a bit.
  \item La IDCT cumple IEEE~1180, y las imágenes intra difieren de la
    referencia solo en lo que el estándar permite.
  \item El decodificador es determinista y se recupera de entradas dañadas
    sin \emph{reset}.
  \item Decodifica holgadamente en tiempo real en CIF, y a 0,5--0,77 veces el
    tiempo real en definición estándar.
  \item El periférico completo consume unos 68~mW estimados.
  \item Los 37 flujos MPEG-2 de conformidad decodifican bit-exactos o con
    la diferencia que la norma admite entre dos IDCT (MAD $\leq 0{,}04$);
    ninguno deriva, después de corregir dos errores del diseño original.
\end{itemize}
